% Copyright 2026 Yilin Liu and GPT-6 Astra Pro. Apache-2.0.
% Curated from sat_research_checkpoint.zip: paper/elementary_crossing_operator_lemma.tex
% Complete finite proof retained; no SAT dependency or formal receipts imported.
\subsection{An elementary proof of the crossing operator identity}
\label{sec:elementary-crossing-operator}

We give a self-contained proof of the finite operator identity used in
the integral crossing construction. This isolates the algebraic
consequence from the satisfiability algorithm: no claim about that
algorithm is needed in this subsection. The identity is the one in
Appendix~B, Lemma~9 of~\cite{anthropic2026}; the proof below uses only
normal ordering and a finite matching expansion.

Let $V=\mathbb Q[p_1,p_2,\ldots]$, graded by $\deg p_j=j$, and let
$V_E$ denote the homogeneous component of weight $E$; set $V_E=0$
for $E<0$. For $k\ge1$ define
\[
 D_{-k}f=p_kf,\qquad D_kf=k\frac{\partial f}{\partial p_k}.
\]
Thus $[D_k,D_{-l}]=k[k=l]\operatorname{Id}$, while creation
operators commute with one another, as do annihilation operators.
For $m\ge1$ write
\begin{align*}
 J_m^+(z)&=\sum_{l\ge1}\binom{l-1}{m-1}D_{-l}z^{l-m},\\
 J_m^-(z)&=\sum_{k\ge1}\binom{-k-1}{m-1}D_kz^{-k-m}.
\end{align*}
The first sum actually starts at $l=m$. For a partition
$\mu=(\mu_1,\ldots,\mu_t)$, define
\[
 Y(p_\mu,z)=:\!\prod_{i=1}^t
       (J_{\mu_i}^+(z)+J_{\mu_i}^-(z))\!:
       =\sum_{r\in\mathbb Z}p_\mu[r]z^{r-|\mu|}.
\]
The colons mean that in each term all creation operators are placed
to the left of all annihilation operators, without adding commutator
terms. Extend the modes linearly to $V$. Every $a[r]$ maps
$V_E$ to $V_{E+r}$. Its action on a polynomial is well-defined:
the positive oscillator indices in a nonzero term have sum at most
the input weight, and the prescribed mode then bounds the negative
indices as well.

\begin{lemma}[Finite crossing operator identity]
\label{lem:elementary-crossing-operator}
Let $h\ge1$, $b\in V_h$, and $a=\sum_{h'=0}^h a_{h'}$ with
$a_{h'}\in V_{h'}$. Set
\[
 a^{\langle i\rangle}=\sum_{h'=0}^h\binom{h+h'}i a_{h'},\qquad
 Q_s(a,b)=\sum_{i=0}^{2h}(-1)^i
    a^{\langle i\rangle}[s-h+i]b[h-s-i].
\]
On every $V_E$ the following statements hold:
\begin{enumerate}
\item Every product in the definition of $Q_s$ can be reversed
simultaneously without changing $Q_s$.
\item $Q_s=0$ if $|s|>E+h$.
\item $\sum_{|s|\le E+h}Q_s=(a[-h]b)\operatorname{Id}_{V_E}$.
\end{enumerate}
\end{lemma}

\begin{proof}
Linearity reduces the claim to $a=p_\mu$ of weight $h'\le h$
and $b=p_\nu$ of weight $h$. Put $H=h+h'$. Let $\iota_+$
mean expansion in nonnegative powers of $w/z$, and $\iota_-$
expansion in nonnegative powers of $z/w$. Directly from the
commutation relations,
\begin{equation}
\label{eq:elementary-contraction}
 [J_m^-(z),J_n^+(w)]
 =\iota_+\frac{\gamma_{m,n}}{(z-w)^{m+n}},\qquad
 \gamma_{m,n}=(-1)^{m-1}
       \frac{(m+n-1)!}{(m-1)!(n-1)!}.
\end{equation}
Indeed, for $m=n=1$ the contraction is
$\sum_{k\ge1}kz^{-k-1}w^{k-1}=\iota_+(z-w)^{-2}$.
Applying the divided derivatives
$\partial_z^{m-1}\partial_w^{n-1}/((m-1)!(n-1)!)$
gives~\eqref{eq:elementary-contraction}. Reversing the fields gives
the expansion $\iota_-$ of the same rational function, since
$(w-z)^{-2}=(z-w)^{-2}$ before differentiation.

Label equal parts of $\mu$ and $\nu$ separately. Commuting
annihilation operators of the first field past creation operators
of the second gives a sum over partial matchings between these two
lists of parts. A matched pair $(m,n)$ contributes the scalar
in~\eqref{eq:elementary-contraction}; all unmatched operators form
one completely normally ordered product. This assertion follows by
repeatedly using $AB=BA+[A,B]$: every commutator is scalar, so it
creates no further noncentral terms, and each labelled operator can
be contracted at most once.

Fix an input and an output monomial of weight $E$. The matrix
coefficient of any completely normally ordered product of the
unmatched operators is a Laurent polynomial in $z,w$. Its
annihilation indices sum to at most $E$. Since input and output
weights agree, its creation indices have that same sum. Hence all
indices are bounded by $E$, and only finitely many terms occur.
The original two ordered products have well-defined matrix
coefficients in the corresponding iterated Laurent-series rings:
the rightmost field has powers bounded below, and fixing one such
power bounds its intermediate weight. The preceding commutations
may therefore also be justified coefficientwise by truncating
oscillator indices and then taking stabilized coefficients.

For a partial matching let $D$ be the sum of the weights of all
matched parts, so $D\le H$. Multiplying its contraction factors
by $(z-w)^H$ leaves $(z-w)^{H-D}$ times an ordinary Laurent
polynomial. Consequently
\[
 (z-w)^H [u]Y(a,z)Y(b,w)v
 \quad\hbox{and}\quad
 (z-w)^H [u]Y(b,w)Y(a,z)v
\]
are the two expansions of one and the same Laurent polynomial
$F_{u,v}(z,w)$. Here $u,v$ are the chosen weight-$E$ monomials,
and $[u]$ extracts a coefficient. This polynomial is homogeneous
of total degree zero, hence has the form
$\sum_s F_s z^s w^{-s}$ with finite support.

Reading the coefficient $z^sw^{-s}$ in the first order gives
\[
 F_s=\sum_{i=0}^H(-1)^i\binom Hi
       [u]a[s-h+i]b[h-s-i]v=[u]Q_s(a,b)v.
\]
Reading it in the other order reverses every product. This proves
the first assertion. In the first order the rightmost mode is
$b[-r]$ with $r=s-h+i$, so it vanishes on $V_E$ when $r>E$.
Thus every term vanishes for $s>E+h$. In the reverse order the
rightmost mode is $a[r]$, which vanishes when $r<-E$. Since
$r\le s+h'\le s+h$, every term vanishes for $s<-E-h$.
This proves the second assertion.

At $z=w=1$, every partial matching with an unmatched part vanishes,
because then $H-D>0$. Fully matched terms leave only a scalar
times the identity. Denote the sum of their scalar coefficients by
$c(a,b)$. It is independent of $E$, so
\[
 \sum_s Q_s(a,b)=c(a,b)\operatorname{Id}_{V_E}.
\]
To identify the scalar, use the vacuum vector $1\in V_0$.
Its two-field matrix coefficient contains only fully matched terms,
and is therefore $\iota_+c(a,b)(z-w)^{-H}$. On the other hand,
$Y(b,w)1=b+O(w)$: each creation summand starts with its own
$p_n$ at power zero. Taking the coefficient of $w^0$ and then
the constant polynomial coefficient gives
$a[-h]b\,z^{-H}=c(a,b)z^{-H}$. Hence $c(a,b)=a[-h]b$,
which proves the third assertion. Summing homogeneous components
of $a$ gives the stated generality.
\end{proof}

\begin{lemma}[The finite binomial duality]
\label{lem:elementary-crossing-duality}
Define the multiplicative linear transform $T:V\to V$ by
\[
 T(p_n)=(-1)^{n-1}\sum_{j=1}^n\binom njp_j.
\]
For partitions $\mu,\nu$ of $h$, let
$z_\mu=\prod_{k\ge1}k^{m_\mu(k)}m_\mu(k)!$. Then
\[
 T(p_\mu)[-h]p_\nu=z_\mu[\mu=\nu].
\]
\end{lemma}

\begin{proof}
This is the unscaled duality in the source's Lemma~10, which we
rederive to make the dependency explicit. Vandermonde's identity,
including negative upper binomial arguments, gives
\[
 \sum_{j=1}^n\binom nj\binom{-k-1}{j-1}
   =\binom{n-k-1}{n-1}.
\]
It is the coefficient of degree $n-1$ in
$(1+t)^n(1+t)^{-k-1}$. The result is zero for $1\le k<n$
and equals $(-1)^{n-1}$ for $k=n$. The sign in $T(p_n)$
therefore makes its surviving coefficient at $k=n$ equal to one.

In a mode of total shift $-h$ acting on $V_h$, any creation
operator would require total differentiation weight greater than
$h$, so it vanishes. In the remaining all-annihilation terms of
$T(p_\mu)[-h]$, each index $k_i$ is at least $\mu_i$, and their
sum is $h=\sum_i\mu_i$. Hence every $k_i=\mu_i$, and
\[
 T(p_\mu)[-h]p_\nu
   =\left(\prod_i D_{\mu_i}\right)p_\nu
   =z_\mu[\mu=\nu].
\]
If all derivative multiplicities fit those of $p_\nu$, equality
of the total weights forces the two partitions to be identical;
otherwise a derivative annihilates the monomial.
\end{proof}

The argument is first carried out over $\mathbb Q$ to define the
displayed operators without changing bases. When an integral lattice
is preserved by their modes, the resulting equality of integral
matrices specializes to every commutative coefficient ring. The
Schur-lattice construction below provides exactly
that lattice and a dual family with scalar $a[-h]b$ equal to a
Kronecker delta.
